% Source: 30_Chapter_21_A_Song_to_Lady_Poverty_The_Apostolic_Li.txt
\bookchapter[A Song to Lady Poverty]{CHAPTER 21}{A SONG TO LADY POVERTY}{The Apostolic Lifestyle}

\leadin{E}{VERYONE, WE KNOW, LOVES} a lover. Songs of endless devotion are not new. The ballads of
the troubadours filled the twelfth century. They were all about courtly love: gallant knights,
sinister magicians, enchanted castles, talking animals, and, always, beautiful damsels.

Perhaps that is why so many felt drawn to St. Francis of Assisi. His song was of Lady Poverty, his
ideal for true Christianity and personal holiness. Early in his preaching career, “Christ’s very
courageous knight,” as his first biographer called him, gave himself devotedly, ardently to Poverty.

According to one of many stories told about him, Francis, with some companions, went out in search
of Poverty. Two old men told Francis that she lived high in the mountains. Scaling the peaks,
Francis found her there “on the throne of her neediness.” She welcomed the travelers and Francis
immediately praised her as “the queen of the virtues.” She replied that she had indeed been with
Adam in Paradise, but after he sinned she had become a homeless wanderer. Then Jesus came. He made
her his elect one, and through her believers multiplied. Monks especially joined her until her enemy
Avarice made them rich and worldly. Then she no longer had a choice; she withdrew from monasticism.
Upon hearing Poverty’s story, Francis vowed to be faithful to her. He took her as his bride, and she
descended the mountain with him.

Francis was not the only spiritual knight devoted to Lady Poverty. In the twelfth and thirteenth
centuries she had a host of admirers. Itinerant preaching and voluntary poverty appealed to the
imaginations and consciences of many Christians. Growing numbers of laymen, instead of relying on
the prayers of monks and bishops, read the Bible in the vernacular and vowed to follow the gospel
mandate: “Sell what you have, give it to the poor, and come follow me.” Some of these believers
chose orthodoxy, others opted for heresy, and at times only a knife edge seemed to separate the two.

One thing was clear: Innocent III’s vision of the ascended Christ who ruled, through his Vicar, all
nations, all learning, and all grace in this life and the next, faced a formidable rival in the
ancient image of the Savior who said, “Foxes have holes and birds of the air have nests but the Son
of man has no place to lay his head.” Where, after all, is true Christianity? In a sacramental
institution or in a self-denying lifestyle?

\section{POVERTY, HERESY, AND VIOLENCE}

The medieval poverty movement is a timeless reminder that political Christianity is only partial
Christianity. The Christian faith is more than papal policy—much more. What shall it profit the
church, as well as men, to gain the world and lose her soul? What advantage are canon laws, holy
crusades, episcopal appointments, and scholastic disputations if common people receive a stone when
asking for bread?

The gospel of voluntary poverty drew its strength from a deep and widespread resentment of a corrupt
and neglectful priesthood. The back-to-the-apostles movement was often allied with political and
economic restlessness in a rapidly changing and expanding society. But at its heart was the
spiritual hunger of people.

At a time when it was desperately needed, pastoral care was a lost art. Robert Grosseteste, the able
bishop of Lincoln, England (1235–1253), decried the covetousness, greed, and immorality of the
clergy: “As the life of pastors is the book of the laity, it is manifest that such as these are the
preachers of all errors and wickedness.” And the foundation of all of this, said Grosseteste, is the
Roman Court. It appoints not pastors but destroyers of men.

Grosseteste’s complaint had a familiar ring to it. As early as the tenth century, monastic reformers
had called for a return to the poverty of the early church. Every zealous preacher knew if apostolic
poverty is the Christian ideal then bishops in their embellished palaces and monks in their wealthy
cloisters were not living the Christian life.

In earlier centuries, however, all such appeals to sacrifice were conveniently channeled into some
new monastic reform. Within the church itself, it was business as usual. The twelfth and thirteenth
centuries proved to be different. Not all preachers of apostolic poverty were willing to remain
within the acceptable boundaries of the church. When they turned against the church they entered the
ranks of heretics.

Modern Christians find it all but impossible to understand the medieval attitude toward heresy. We
believe deeply that religious faith is a matter of personal choice but we seldom think of religious
beliefs as life-and-death matters. Why should anyone either die for his own faith or kill another
for his?

Medieval Christians, however, never considered faith strictly private. Christian beliefs were the
cement of society; denying a doctrine was tantamount to treason. Christendom was, to shift the
image, a sociopolitical body, and the Christian faith was its life-giving soul. So heresy in
Christendom was no more acceptable than cancer in the flesh.

But what is heresy? In the twelfth century, it was the denial by a baptized person of any revealed
truth of the Christian faith. Among these truths were the unity of the church and the divine
appointment of the pope as head of the Christian body. Therefore, disobedience to established
authority was itself heresy.

In dealing with heretics, then, the church had two primary objectives: the conversion of the heretic
and the protection of Christian society. But how far can the church go to protect society? Is it
right to take a life in order to protect other lives?

Heresy drove the Catholic Church to her most serious internal conflict: how can the church employ
violence to safeguard a peaceful society? The church deliberately accepted a line of action all but
impossible to reconcile with the eternal kingdom toward which she aspired. She created the
Inquisition, not only to execute heretics but to subject them to deliberate and prolonged torture.
In driving out one devil the church opened the door for seven others.

The contradiction was not widely apparent at the time. The same church that sent crusading armies
against the infidels could command the burning of heretics. Almost everyone agreed that a pure
church was the will of God.

\section{THE POOR IN THE KINGDOM}

But where is the one pure church? Is it in the papal palace in Rome? Is it in the blood of crusading
armies or in the sale of indulgences to the poor? What if the one pure church is in none of these,
but is in fact in the hungry who find bread, the naked who are clothed, and the stranger who finds
rest? What if the kingdom of heaven does belong to the poor in spirit? That was the question Lady
Poverty raised.

One of the earliest voices against the worldliness of the Catholic Church was Arnold, an abbot at
Brescia, a town in northern Italy. In a series of sermons at Brescia, Arnold insisted that clerical
vice was a result of the church’s attempt to control the world. He urged the church to surrender its
property and secular dominion to the state and return to the poverty and simplicity of the early
church. The true church and its ministers, he said, should shun wealth, for wealth and power nullify
salvation.

By 1139 Arnold succeeded in turning the people against their bishop. For this, the pope, Innocent
II, banished Arnold from Italy. He apparently fled to Paris where he studied under Abelard and
aroused the wrath of Bernard of Clairvaux, just as his teacher had succeeded in doing. Another
“ravening wolf in sheep’s clothing” Bernard called him.

After five years in exile Arnold showed up in Rome and immediately joined a movement to overthrow
papal dominion. The Romans, filled with dreams of the ancient Roman republic, seized power during
the pope’s preaching tour for the Second Crusade and Arnold stepped into leadership of the new
purely secular government. He preached that clergy should live in apostolic poverty, and he
denounced the College of Cardinals as a den of thieves.

The experiment lasted about ten years before Pope Hadrian IV placed Rome under an interdict and
secured the help of Emperor Frederick Barbarossa in capturing Arnold. In 1155 he was executed by
burning and his ashes were thrown into the Tiber River.

People had scarcely forgotten Arnold when another voice for poverty arose in eastern France: Peter
Waldo (c. 1140–1218), a rich merchant of Lyons.

One day Waldo heard a wandering troubadour singing the virtues of the monastic life. The ballad was
about young Alexis whose patrician Roman parents pressed him into marriage. The reluctant groom,
however, was dedicated to the ideal of chastity, so on his wedding night he made a pact of virgin
purity with his bride and immediately left for the Holy Land.

Alexis’ parents searched for him in vain. Years later he returned home a beggar, so emaciated from
his life of self-denial that no one recognized him. He existed in the courtyard on scraps from the
family table.

Only as he lay dying did he reveal his true identity, too late for the grieving family to claim him.
The moral was pointed: a true Christian must be willing to sacrifice everything in this life for the
sake of the next.

Struck to the heart by the story, Waldo sought a priest to find out how to live like Christ. The
priest turned him to the answer Jesus gave to the rich young ruler: “If you would be perfect, go,
sell what you possess and give to the poor, and you will have treasure in heaven; and come, follow
me” (Matt. 19:21). The same text nine centuries before had launched the monastic movement with
Anthony in Egypt.

Waldo decided to follow the counsel. He provided an adequate income for his wife, placed his two
daughters in a cloister, and gave the rest of his estate to the poor.

To launch his mission to the poor, Waldo enlisted two priests to translate portions of the Bible
into French. After memorizing long passages, Waldo began to teach common folk to imitate Christ by
practicing voluntary poverty. Thus his innovations lay in applying the life of poverty,
discipleship, and preaching to all true Christians, not just monks.

As he gained a few followers, Waldo sent them out two by two, after the apostolic pattern, into
villages and marketplaces, to teach and explain the Scriptures. They called themselves the “Poor in
Spirit.” We know them as Waldenses.

Waldo’s unauthorized preaching soon met the stiff opposition of the Archbishop of Lyon, who ordered
him to stop. Waldo refused, quoting St. Peter: “We must obey God rather than men” (Acts 5:29). The
Archbishop proceeded to excommunicate him.

Waldo and his followers decided to appeal to the pope. They arrived in Rome and found it crowded
with churchmen attending the Third Lateran Council (1179). They were able to gain a hearing before
the Council but had the misfortune to be ridiculed by a smooth, fast-talking Englishman named Walter
Map.

Pope Alexander III found no evidence of heresy among them and was impressed by their poverty. They
were mere laymen, however, so he ruled that they could preach only by the invitation of bishops, a
very unlikely prospect.

Waldo was convinced that the Scriptures commanded him to preach to the poor with or without bishops’
approval. Along with a growing number of followers, he continued to preach and practice apostolic
poverty. The movement spread into southern France and across the Alps into Italy. By 1184 their
disobedience compelled Pope Lucius III to excommunicate them from the holy Catholic Church.

The conflict is understandable. The Waldenses wanted to purify the church by a return to the simple
life of the apostles. This meant the surrender of worldly power. Their aim, like that of the Roman
church, was salvation. But their means were radically different. The papacy could not renounce its
sacraments or its priesthood, nor admit that faith in God might be something other than the mandates
of Rome. From their side, the Waldenses came to feel more and more that no teaching except Christ’s
was binding. The Scriptures must rule. But how could they find support for their cause if everyone
lived in apostolic poverty? Slowly they came to accept, just as early monastic houses had, two
levels of Christian commitment. The “Poor in Spirit,” the society proper, was bound by special vows
and worshiped together in simple services. Another circle of “friends” remained in the Catholic
Church but supplied recruits and support for the movement.

The Waldenses were so clearly a back-to-the-Bible movement that over the years many evangelical
Christians have tried to present them as “reformers before the Reformation.” Compared to the Roman
church’s doctrine of papal authority, the Waldensian call to return to the Bible does indeed sound
like Luther or Calvin. But their view of salvation, a life of penance and poverty, lacks the clear
note of God’s grace that sounded so powerfully in the Reformation.

\section{THE PERIL OF NEW HERESIES}

A third dissenting movement, and by far the most troublesome for the Catholic Church, was the group
commonly called Cathari, which means pure ones. Since they were especially influential in and about
the town of Albi in southern France, some people called them Albigenses.

Most of what we know about the Albigenses comes through their enemies and, as Charles Williams once
observed, not one mind in a thousand can be trusted to state accurately what one’s opponent says,
much less what he thinks.

Very likely, however, the Cathari infiltrated Europe from Bulgaria, where one of their chief
branches bore the name Bogomiles. Like the Gnostics in the early church, the Cathari held that the
universe is the scene of an eternal conflict between two powers, the one good, the other evil.
Matter, including the human body, is the work of this evil power, the god of the Old Testament. He
had, they claimed, imprisoned the human soul in its earthly body.

To escape from the power of the flesh, the true Cathar was supposed to avoid marriage, sexual
intercourse, eating of meat, and material possessions. Here was a radical poverty, but not one based
on the example of Jesus so much or the nature of the universe. A cosmic civil war raged between
matter and spirit, and the Cathari joined the ranks of spirit.

The good God, they taught, sent Christ to reveal the way of salvation to man; Christ, to a Cathar,
was not a human being but a life-giving spirit. A human Christ was an impossibility, and salvation
by death on the cross was unthinkable. Christ taught the way, like a Buddha rather than the God-man
of Christian creeds.

Obviously the Cathari were heretical in a way Arnold of Brescia and Peter Waldo were not. Arnold and
Waldo refused to submit to church authorities, but the Cathari rejected not only popes and bishops
but basic Christianity. They tried to escape from evil, not by repentance and faith but by
discovering and liberating the pure self.

The Cathari were an immense peril to the Roman church. Not only had they revived the ancient dualist
heresy; by 1200 they had gained the protection of the princes of Toulouse, a cultural area in
southern France, and were spreading at an alarming rate. Three weapons were at the Catholic Church’s
disposal: preaching to return them to the truth, a crusade to crush all hardened resistance, and the
Inquisition to uproot heresy completely.

\section{RISE OF THE DOMINICANS}

The popes tried sending preachers among the Cathari, but they were notoriously ineffective, at least
until a Spaniard named Dominic Guzman (1170–1221) struck upon the reason. In 1206 Dominic was
helping in efforts to convert the Albigenses when he realized that the papally assigned preachers
were depending upon their ecclesiastical pomp and dignity. The Albigenses considered such show a
sure sign of false religion. Dominic believed that the heretics would listen if the preachers
themselves were committed to poverty. To win the heretics, Dominic went forth among them as a poor
man, barefoot and begging.

Dominic’s peaceful mission lasted only two years in southern France before it was forced aside by
Innocent III’s get tough policy. However, the zealous Spaniard was convinced that poverty and
preaching belonged together. He gathered a group of like-minded men and continued his work among the
heretics in other places. In 1220 the Dominican mission and lifestyle gained official approval. The
new preaching order that we know as Dominicans was called mendicant, meaning begging, and the term
friar (or brother) distinguished them from monks because, unlike monks, they went forth to live
among people to preach and teach. Just as monastic houses had once arisen to minister in the
countryside, so the mendicant friars now emerged to meet the spiritual needs of townspeople.

Meanwhile, Innocent III was determined to crush the Albigensian heresy. The northern French were
itching for a chance to gobble up southern France, which was then a separate country. When Innocent
called for a crusade—not against the Muslim Turks but against Christian heretics—the northern French
came pillaging and murdering. It isn’t often that good Christian men can save their souls and at the
same time enlarge their kingdoms by butchering unbelievers. Even Innocent was shocked by the
brutality. But the crusade was eminently successful. By 1215 the Albigenses were uprooted in
Toulouse and the northern French had claimed the ravaged territories of the south.

\section{ENTER THE INQUISITOR}

The scattered heretics had to be hounded out of hiding. That was the purpose of the Inquisition. The
infamy of this institution has left its mark on the memories and vocabularies of men everywhere. We
equate it with the ruthless miscarriage of justice.

The early form of the Inquisition appeared in 1184 when Pope Lucius III required bishops to
“inquire” into the beliefs of their subjects. In short, they held an inquiry or inquest. Heresy or
harboring a heretic brought immediate excommunication.

The spread of the Albigenses and their own violence toward loyal Catholics called for stricter
measures. In 1215 the Fourth Lateran Council, under Innocent III’s leadership, provided for the
state’s punishment of heretics, the confiscation of their property, excommunication for those
unwilling to move against the heretic, and complete forgiveness of sins for those cooperating.

In 1220 the pope took the Inquisition from the hands of the bishops and turned it over to the
newly-formed Dominicans, and nine years later the Synod of Toulouse systematized inquisitorial
policies, leaving the alleged heretic with virtually no rights. The inquisitor was subject to no
law, only to the pope. He was prosecutor and judge. The “trial” was secret, and the accused had to
prove his innocence, as in all courts following Roman law, without the benefit of counsel or
knowledge of his accusers.

The final significant step came in 1252. Pope Innocent IV authorized torture as a means of getting
information and confessions from accused heretics. Popes, saints, and theologians had in the past
rejected with horror the very thought. But no such reserve remained after Innocent III had ascended
the papal throne and the Catholic Church had achieved its majestic and powerful unity.

Canon Law, it is true, forbade a cleric from shedding blood. He who served the altars of the One
Sacrifice must not sacrifice men. He could only hound and interrogate and torture the prisoner. If
he found the unfortunate person guilty of heresy, he turned him over to civil authorities, usually
for burning at the stake.

It was an ugly business, but almost everyone, after Augustine, agreed that saving the body by
amputating a rotten limb was the path of wisdom. Clearly the Church of Rome was the body and the
heretic the rotten limb.

The combination of the crusade against heresy in Toulouse and the Inquisition brought an end to
Catharism before the thirteenth century closed. The doctrine that matter was incapable of salvation
was, until the coming of Mary Baker Eddy, a thing of the past. The Waldenses, who were often targets
of the same Inquisition, survived in the mountains of Italy and welcomed the Reformation when it
broke over Europe in the sixteenth century.

The Inquisition also survived, even with its serious weakness. It could amputate, but it could not
heal. The ministry of healing arose in a village surrounded by vineyards eighty-five miles north of
Rome. Assisi was the hometown of Giovanni Bernardone; we know him as St. Francis of Assisi.

Francis (1182–1226) was the son of a wealthy Italian cloth merchant. Francis’ father had dreams of
knighthood for his son. Francis pursued the glory of military service, but his capture, illness, and
a series of visions led him to turn his back on his father’s dreams. The re-orientation was
dramatic. Before his turning he benefited from wealth, sought honor, and was repulsed by illness
like leprosy. With Christ as his model, he abandoned the security of wealth, lived a humble life
dedicated to benefiting others, and embraced lepers. Leaving home in a ragged cloak and a rope-belt
taken from a scarecrow, he wandered the countryside with a few followers, begging from the rich,
giving to the poor, and preaching the joys of apostolic poverty.

\section{FOUNDING THE FRANCISCANS}

In 1209 Francis prepared a simple Rule for his little brotherhood. It consisted chiefly of Christ’s
appeal to take up the cross, his advice to the rich young ruler, and his directions on sending out
the apostles. With this Rule, Francis and his companions went to Pope Innocent III for approval. The
scene was almost a reproduction of Waldo before Alexander III in 1179, but times had changed.
Innocent saw the mistake others had made, so he approved the little preaching band. Francis called
his group the Friars Minor (Lesser Brothers). We call them the Franciscans.

Almost from the start Francis’ vision was for the world. He tried to go to Syria and to Morocco but
was thwarted by misfortune. Then in 1219 a crusading expedition to Egypt gave him his opportunity.
With eleven companions he accompanied the army to the Middle East where he tried unsuccessfully to
convert the Sultan of Egypt. From Egypt he visited the holy places in Palestine, and it was more
than a year before he again saw Italy.

During his absence differences arose among the brothers. Some felt that their rapid growth demanded
more organization, more rules, more supervision. Others tried to cling to the original principles of
Christ-like poverty.

Upon Francis’ return to Italy, he saw that the problems were more than he could handle. He was a
model, not a manager. So he appealed to the pope to appoint Cardinal Ugolino as his adviser and soon
he surrendered the administration of the brotherhood to his associate Peter de Cataneo. “Lord
Jesus,” he prayed, “I give Thee back this family which Thou didst entrust to me. Thou knowest that I
have no longer the strength or the ability to take care of it.” It is a familiar scene in Christian
organizations. One man establishes; another man administers.

Ugolino, who later became Pope Gregory IX, admired Francis greatly, but he was above all a prince of
the church. He saw the possibilities of the movement as an agent for the advancement of the Roman
church, especially where its authority had been undermined by Arnold, Waldo, and the Cathari. He
would reform the church by giving the Franciscans authority; Francis wanted to reform the world by
preaching Christ-like humility.

In 1223 Pope Honorius III confirmed a new Rule for the order, which allowed for an elaborate
organization and made begging a fundamental trait of the order. Thus Franciscans joined Dominicans
in a great mendicant counteroffensive in preaching, doctrine, and dedication. As his end drew near,
he was carried back to Assisi and died there on October 3, 1226, in full humility and song.

Francis was a product of his age, a lover of Lady Poverty, but he belongs to all the Christian
centuries. “For a few years,” wrote Herbert Workman, “the Sermon on the Mount became a realized
fact. But the dream passed” from all but a few.

For another century the Franciscans were divided between the Conventuals, who allowed the church to
hold property for Franciscan use, and the Spirituals, who argued that complete poverty was an
indispensable note of the true church.

Thus Lady Poverty remained a challenge to the church of wealth and power, but like Innocent III’s
vision of universal dominion, it proved an impossibility for the whole church for all times.

\readinghead{Suggestions for Further Reading}
\begin{readinglist}
\entry Baldwin, Marshall W. The Mediaeval Church. Ithaca, NY: Cornell University Press, 1953.
\entry Gobry, Ivan. St. Francis Assisi. San Francisco: Ignatius, 2006.
\entry Lambert, Malcolm. Medieval Heresy. New York: Holmes \& Meier, 1976.
\entry Runciman, Steven. The Medieval Manichee. London: Cambridge University Press, 1955.
\entry Thompson, Augustine. Francis of Assisi: A New Biography. New York: Cornell University Press,
  2012.
\entry Turberville, A. S. Medieval Heresy and the Inquisition. London: Archon, 1964.
\entry Westin, Gunnar. The Free Church through the Ages. Nashville: Broadman, 1958.
\entry Workman, Herbert B. The Evolution of the Monastic Ideal. Boston: Beacon Press, 1962.
\end{readinglist}
